% tbsrc/regularization.tex — Regularized Reconstruction: MAP-EM
% Status: written (first draft). Follows the "Beyond early stopping" note in
% acceleration.tex. MAP objective f + beta R; De Pierro surrogate -> closed-form
% penalized_mlem for quadratic priors (intensity, smoothness); OSL -> osl_mlem for
% edge-preserving (Huber). Matches src/penalized.jl; De Pierro derivation in App.B.
\section{Regularized Reconstruction: MAP-EM}
\label{sec:regularization}

Early stopping (Section~\ref{sec:acceleration:reg}) controls noise by halting the
unregularized iteration before it fits the Poisson realization; it is blunt
because one global iteration budget governs every region irrespective of its
local statistics. The explicit alternative is \emph{maximum a posteriori} (MAP)
reconstruction: minimize the penalized objective
\begin{equation}
  \Phi(\xv) \;=\; f(\xv) \;+\; \beta\,R(\xv),
  \label{eq:reg-map}
\end{equation}
where $f$ is the negative Poisson log-likelihood of \eqref{eq:st-listmode}, $R$ is
a prior penalizing roughness, and $\beta\ge0$ sets its strength ($\beta=0$
recovers MLEM). The preconditioned-gradient view of Section~\ref{sec:mlem} makes
this a small change in principle --- the prior's gradient simply adds to
$\nabla f$ --- but the multiplicative update earns its non-negativity and
monotonicity from the EM structure, and naively bolting a gradient onto it forfeits
both. This section gives the two updates RecoCrysp implements in
\code{penalized.jl} that keep what the EM structure provides: a closed-form,
monotone update for \emph{quadratic} priors (De Pierro), and a One-Step-Late update
for \emph{edge-preserving} priors (Huber).

\subsection{Quadratic priors and the De Pierro update}
\label{sec:reg:depierro}

The obstacle is that $R$ couples neighbouring voxels, so its gradient mixes them
and the per-voxel decoupling that gave MLEM its closed form (Appendix~\ref{app:em})
is lost. De Pierro's surrogate~\cite{depierro1995} restores it: for a
\emph{quadratic} $R$, a convexity argument (Appendix~\ref{app:depierro}) replaces
the coupled penalty by a separable upper bound that is exact at the current
iterate, so the M-step again solves one independent problem per voxel --- now a
quadratic rather than a linear equation. Writing the ordinary EM numerator
\begin{equation}
  \\bm{t} \;=\; \xv_k \,\odot\, \Aadj\!\Bigl(\code{mult}\odot\frac{\code{counts}}{\pred_k}\Bigr)
  \label{eq:reg-numerator}
\end{equation}
(the quantity MLEM divides by $\sensv$), the penalized M-step is the positive root
of the per-voxel quadratic $G\,x^2 + B\,x - t = 0$,
\begin{equation}
  \xv_{k+1} \;=\; \frac{2\,\\bm{t}}{\sqrt{B^2 + 4\,G\,\\bm{t}}\;+\;B},
  \label{eq:reg-depierro}
\end{equation}
where the prior enters \emph{only} through the coefficients $(B,G)$, listed in
Table~\ref{tab:reg-priors}. The form makes the reduction to MLEM transparent: with
$\beta=0$ the prior contributes $G=0$ and $B=\sensv$, and \eqref{eq:reg-depierro}
collapses to $\xv_{k+1}=\\bm{t}/\sensv$, exactly \eqref{eq:ml-em}. The root is taken
in the rationalized form $2t/(\sqrt{B^2+4Gt}+B)$ rather than
$(-B+\sqrt{B^2+4Gt})/(2G)$ so that it stays well-defined as $G\to0$.

\begin{table}[t]
  \centering
  \begin{tabular}{@{}lll@{}}
    \toprule
    Prior $R(\xv)$ & $B$ & $G$ \\
    \midrule
    none (MLEM)
      & $\sensv$ & $0$ \\[2pt]
    intensity $\tfrac{\beta}{2}\norm{\xv-\bm z}^2$
      & $\sensv-\beta\,\bm z$ & $\beta$ \\[2pt]
    smoothness $\tfrac{\beta}{2}\sum_j\sum_{k\in\mathcal N_j}(x_j-x_k)^2$
      & $\sensv-\beta\,(\bm w\odot\xv_k+\bm m)$ & $2\beta\,\bm w$ \\
    \bottomrule
  \end{tabular}
  \caption{De Pierro coefficients $(B,G)$ for \eqref{eq:reg-depierro}, one entry
    per prior in \code{penalized.jl}. For the smoothness prior $\mathcal N_j$ is
    the 6-connected (face) neighbourhood, $\bm m_j=\sum_{k\in\mathcal N_j}x_k$ the
    neighbour sum, and $\bm w_j=\abs{\mathcal N_j}\in\{3,\dots,6\}$ the in-bounds
    neighbour count, both evaluated at the current iterate $\xv_k$. Every row
    reduces to $(\sensv,0)\Rightarrow$ MLEM at $\beta=0$.}
  \label{tab:reg-priors}
\end{table}

Two quadratic priors are provided. The \code{QuadraticIntensityPrior}$(\beta,\bm
z)$ pulls the image toward a reference $\bm z$; it is De Pierro's example prior and
serves mainly as a baseline, since penalizing intensity biases the activity. The
\code{QuadraticSmoothnessPrior}$(\beta)$ is the noise-controlling one: it penalizes
squared differences across the 6-connected neighbourhood, so its gradient is a
discrete Laplacian and the prior acts as a smoother. Because the De Pierro update
is the exact M-step of a surrogate that majorizes $\Phi$ and touches it at $\xv_k$,
it is \emph{monotone} --- $\Phi(\xv_{k+1})\le\Phi(\xv_k)$ at every iteration ---
inheriting the convergence guarantee MLEM has for the likelihood alone.

The neighbour sums $\bm m$ and counts $\bm w$ are the only new computation, and
they are a stencil over the image --- the same kind of structured, embarrassingly
parallel operation as the Joseph projector. They are written as a
KernelAbstractions kernel so the one source runs on CPU or GPU
(Section~\ref{sec:acceleration}), mirroring the projector idiom:
\begin{lstlisting}
@kernel function _neigh6_kernel!(nsum, wsum, @Const(x), n)
    I = @index(Global, Cartesian)
    i, j, k = I[1], I[2], I[3]
    s = 0.0f0; w = 0.0f0
    if i > 1;    s += x[i-1,j,k]; w += 1.0f0; end   # ... 5 more faces
    if i < n[1]; s += x[i+1,j,k]; w += 1.0f0; end
    nsum[I] = s; wsum[I] = w
end
\end{lstlisting}
The De Pierro iteration is then \eqref{eq:reg-numerator}--\eqref{eq:reg-depierro}
verbatim, with the prior reduced to its $(B,G)$ pair:
\begin{lstlisting}
B, G  = _prior_BG(prior, x, sens)               # Table 1; NoPrior -> (sens, 0)
denom = sqrt.(B.^2 .+ 4f0 .* G .* t) .+ B
x     = ifelse.(sens .> 0f0, 2f0 .* t ./ denom, x)   # leave FOV-exterior fixed
x     = max.(x, 0f0)                                  # guard Float32 cancellation
\end{lstlisting}
This is \code{penalized\_mlem}; called with \code{NoPrior()} it is bit-for-bit the
MLEM of Section~\ref{sec:mlem}.

\subsection{Edge-preserving priors: One-Step-Late}
\label{sec:reg:osl}

A quadratic smoothness prior penalizes \emph{every} difference equally, so it
blurs genuine edges along with noise. Edge-preserving priors penalize small
differences quadratically but large ones only linearly, leaving true boundaries
intact; the Huber penalty $\rho_\delta$ --- quadratic for $\abs t\le\delta$, linear
beyond --- is the canonical choice. Such priors are not quadratic, so the De Pierro
closed form does not apply. RecoCrysp uses instead the One-Step-Late (OSL) update
of Green~\cite{green1990}, which evaluates the penalty gradient at the
\emph{current} iterate and adds it to the sensitivity in the MLEM denominator:
\begin{equation}
  \xv_{k+1} \;=\; \frac{\\bm{t}}{\sensv + \nabla R(\xv_k)},
  \qquad
  \bigl(\nabla R\bigr)_j = \beta\!\!\sum_{k\in\mathcal N_j}\!\psi_\delta(x_j-x_k),
  \quad
  \psi_\delta(t)=\operatorname{clamp}(t,-\delta,\delta).
  \label{eq:reg-osl}
\end{equation}
The clamped influence $\psi_\delta$ is the Huber gradient: below $\delta$ it is the
identity, so the prior smooths like the quadratic one; above $\delta$ --- at an
edge --- it saturates, so the penalty stops pulling across the boundary and the
edge survives. The gradient is the same 6-neighbour stencil as before with
\code{clamp} in place of the bare difference, and $\beta=0$ again recovers MLEM.
OSL is the standard MLEM-structured edge-preserving scheme, but because the penalty
gradient is taken one step late it is \emph{not} provably monotone: it is stable at
the moderate $\beta$ used in practice and can diverge at very large $\beta$, so we
lean on $\beta=0\equiv$ MLEM as the anchor and verify behaviour empirically. For
quadratic priors \code{penalized\_mlem} (monotone) is preferred; \code{osl\_mlem}
is for the non-quadratic, edge-preserving case.

\subsection{Choosing the strength}
\label{sec:reg:beta}

The regularization weight $\beta$ has no universal value: it trades resolution
against noise, and because the prior is added to $f$ it must be commensurate with
the magnitude of $\sensv$, which depends on the geometry, the multiplicative
factors, and the event-to-sensitivity scale (Section~\ref{sec:datamodel}). On the
worked NEMA study (Part~II) $\sensv$ is of order $10^2$--$10^3$, so useful
smoothness/Huber strengths fall in $\beta\sim\mathcal O(10^2$--$10^3)$ and the Huber
threshold $\delta$ is set just above the noise scale of the differences;
both are tuned per study, like the iteration count. The practical finding there is
that for low-contrast features a quadratic smoothness prior, an edge-preserving
Huber prior, and a simple Gaussian post-filter on an early-stopped MLEM all sit on
nearly the same resolution--noise frontier --- regularization \emph{bounds} the
noise growth that early stopping merely outruns, and which of these realizations is
used matters little. Edge preservation pays off only where the features carry
strong edges for the Huber threshold to protect. The detailed comparison is the
subject of Part~II.
